% C4 Architecture Visualization Documentation Suite | Original educational content
% Standalone LaTeX root: no external style file, no network, no shell escape.
\documentclass[10pt,letterpaper]{article}
\usepackage[T1]{fontenc}
\usepackage[utf8]{inputenc}
\usepackage{lmodern,microtype}
\usepackage[letterpaper,margin=0.82in,headheight=15pt,headsep=18pt,footskip=26pt]{geometry}
\usepackage[table]{xcolor}
\usepackage{booktabs,tabularx,array,longtable}
\usepackage{enumitem}
\usepackage{fancyhdr,lastpage}
\usepackage{titlesec}
\usepackage{hyperref,xurl}
\usepackage{needspace}
\usepackage{listings}
\usepackage{amssymb}
\definecolor{C4Navy}{HTML}{17364B}
\definecolor{C4Teal}{HTML}{137C85}
\definecolor{C4Gray}{HTML}{526775}
\definecolor{C4Light}{HTML}{EEF4F5}
\definecolor{C4Line}{HTML}{CAD8DE}
\color{C4Navy}
\hypersetup{colorlinks=true,linkcolor=C4Teal,urlcolor=C4Teal,citecolor=C4Teal}
\setlength{\parindent}{0pt}
\setlength{\parskip}{5pt}
\setlength{\emergencystretch}{2em}
\raggedbottom
\setlist[itemize]{leftmargin=1.4em,topsep=3pt,itemsep=3pt}
\setlist[enumerate]{leftmargin=1.6em,topsep=3pt,itemsep=3pt}
\titleformat{\section}{\large\bfseries\sffamily\color{C4Navy}}{\thesection}{0.6em}{}
\titleformat{\subsection}{\normalsize\bfseries\sffamily\color{C4Teal}}{\thesubsection}{0.6em}{}
\titlespacing*{\section}{0pt}{16pt}{6pt}
\titlespacing*{\subsection}{0pt}{9pt}{4pt}
\pagestyle{fancy}
\fancyhf{}
\fancyhead[L]{\scriptsize\sffamily C4 / ARCHITECTURE VISUALIZATION}
\fancyhead[R]{\scriptsize\sffamily REFERENCE / EDUCATION}
\fancyfoot[L]{\scriptsize\color{C4Gray}ORIGINAL ENGINEERING SUITE / 2026}
\fancyfoot[R]{\scriptsize\thepage\ / \pageref*{LastPage}}
\renewcommand{\headrulewidth}{0.4pt}
\newcolumntype{Y}{>{\raggedright\arraybackslash}X}
\newcommand{\RuleLine}{\par\noindent\color{C4Line}\rule{\linewidth}{0.45pt}\color{C4Navy}\par}
\newcommand{\WritingSpace}[1]{\par\vspace{2pt}\noindent\fcolorbox{C4Line}{white}{\parbox[t][#1][t]{0.94\linewidth}{\hspace{1pt}}}\par\vspace{7pt}}
\newcommand{\Field}[1]{\par\textbf{#1}\par\WritingSpace{0.57in}}
\newcommand{\Prompter}[1]{\par{\small\color{C4Teal}\textbf{Authoring guidance:} #1}\par}
\newcommand{\SourceNote}[1]{\par\smallskip\noindent\colorbox{C4Light}{\parbox{0.94\linewidth}{\footnotesize\textbf{PROVENANCE.} #1}}\par\medskip}
\lstset{basicstyle=\ttfamily\footnotesize,breaklines=true,columns=fullflexible,showstringspaces=false,frame=single,rulecolor=\color{C4Line}}
\newcommand{\DocStart}[3]{%
\begin{document}\begin{titlepage}
\vspace*{0.55in}
{\sffamily\bfseries\small\color{C4Teal}#1}\par\vspace{0.24in}
{\sffamily\bfseries\fontsize{26}{31}\selectfont\color{C4Navy}#2}\par\vspace{0.18in}
{\large\color{C4Gray}#3}\par\vspace{0.4in}
\textcolor{C4Teal}{\rule{\linewidth}{1.15pt}}\par\vspace{0.24in}
\textbf{Document control}\par
Edition 1.0\quad|\quad 30 October 2026\quad|\quad Original educational engineering material\par
\vspace{0.17in}
\textbf{Source boundary}\par
Simon Brown, \emph{The C4 Model: Visualizing Software Architecture} (O'Reilly, 2026), is the primary study source. The suite's exercises, forms, crosswalks, process rules, sample systems, and scoring models are independently authored supplements; they are not represented as the author's prescriptions.\par
\vfill
{\small\color{C4Gray}Not an ISO/IEC/IEEE 42010 conformance certificate. Not a substitute for the published standard or for review of deployed systems.}\par
\end{titlepage}
}
\DocStart{BOOK-ALIGNED STUDY GUIDE}{Chapter 10: Notation}{Primary reading: book chapter 10, printed pages 97–113}
\SourceNote{The chapter title, structure, and focus below are aligned to Brown (2026). Commentary, tasks, and answer criteria are new instructional material; they do not reproduce book chapters.}
\section{Reading objectives}
\begin{itemize}
\item Explain the purpose and important boundaries of notation in a software-architecture conversation.
\item Recognize the chapter-specific failure modes and explain their likely effect on a consumer of the diagram.
\item Produce original architecture evidence and critique it independently of a preferred drawing tool.
\end{itemize}

\section{Chapter reading map}
\begin{itemize}
\item Read and outline the book section: What Makes a Good Notation?.
\item Read and outline the book section: Diagram Titles.
\item Read and outline the book section: Acronyms, Abbreviations, and Code Names.
\item Read and outline the book section: Layout.
\item Read and outline the book section: Elements.
\item Read and outline the book section: Relationships.
\item Read and outline the book section: Diagram Key.
\item Read and outline the book section: Do We Need a Standard Notation?.
\end{itemize}

\section{Concept synthesis}
Notation matters when it removes interpretation burden: titles, types, directions, descriptions, legends, and treatment of sizes/colors should carry intentional meaning.

\begin{itemize}
\item A consistent notation can still be ambiguous unless element abstractions are sound.
\item Color cannot be the sole carrier of critical distinctions.
\item Relationship direction should match verbs, not merely aesthetic layout.
\end{itemize}

\section{Critical interpretation}
The following questions deliberately test whether the reader understands the abstraction and the stakeholder concern rather than memorizing diagram symbols.

\begin{enumerate}
\item What key architectural question can the techniques from Chapter 10 answer, and which concerns still require another view?
\item What evidence would show a diagram is truthful today rather than merely attractive?
\item What is the smallest level of detail that would answer the intended stakeholder concern?
\item Which element or edge would be most expensive to misinterpret? Explain why.
\end{enumerate}

\section{Practical exercises (original)}
\begin{enumerate}
\item Audit a multi-diagram set for naming and notation drift.
\item Create a monochrome-reviewable diagram key.
\item Run a blind interpretation exercise with two reviewers.
\end{enumerate}

\section{Model inspection protocol}
\begin{itemize}
\item State the intended audience and decision before starting the drawing.
\item Write the view scope, abstraction level, owned software boundary, and out-of-scope items.
\item Inspect identifiers, labeled relationships, technology choices where significant, and legend consistency.
\item Ask a second reviewer to interpret the picture without its author narrating.
\item Record differences between the intended and observed interpretation; revise the source model.
\end{itemize}

\section{Self-assessment and evidence}
\begin{tabularx}{\linewidth}{@{}YY@{}}\toprule
\textbf{Criterion} & \textbf{Passing evidence} \\ \midrule
Conceptual command & A concise explanation of why this view exists and what it intentionally excludes. \\
Semantic precision & Stable names, correctly bounded elements, verb-labeled relationships and no conflated abstraction. \\
Repository evidence & Versioned diagram source, review findings and link to the relevant system inventory. \\
Extension discipline & The student labels their lab policy or new tool convention as supplementary practice. \\
\bottomrule\end{tabularx}


\section{Extension bridge}
Supplementary engineering practice: connect this chapter to stakeholder concerns (ISO/IEC/IEEE 42010), to module/component-and-connector/allocation concerns (Views and Beyond), and to the repository PlantUML style framework. These crosswalks are interpretive mappings; they are not rules stated in the book.


\end{document}
